% amsthm must come before alex.sty (which calls \theoremstyle)
\usepackage{amsthm}

% Index. imakeidx must precede hyperref so index page numbers become links.
% The Makefile runs makeindex explicitly (noautomatic), so the build does not
% depend on shell escape.
\usepackage{imakeidx}
\makeindex[intoc,columns=2,noautomatic,options={-s book.ist}]

% Figures (TikZ; simple line drawings only, no external files).
\usepackage{tikz}
\usetikzlibrary{arrows.meta,positioning,calc,shapes.geometric}
\usepackage{pgfplots}
\pgfplotsset{compat=1.17}

% Hyperlinks: must come before alex.sty because alex.sty loads cleveref,
% and cleveref must be loaded after hyperref.
\usepackage[hidelinks,colorlinks=false]{hyperref}
\hypersetup{
  pdftitle={Bernoulli Types: Two-Sided Error in Approximate Membership, Retrieval, and Computation},
  pdfauthor={Alexander Towell},
  pdfsubject={Research monograph: two-sided error in approximate membership, retrieval, and computation},
  pdfkeywords={approximate membership, two-sided error, false positive rate, false negative rate,
    Bloom filter, retrieval, perfect hash filter, binary channel, randomized response,
    error propagation, approximate types},
  pdfcreator={pdfLaTeX},
  pdflang={en-US}
}

% Notation package (unified macros: ASet, Set, fprate, fnrate, etc.)
% alex.sty loads algorithm2e and cleveref internally.
% Do not also load algorithm/algpseudocode or a second cleveref.
\usepackage[fancy,section]{alex}

% Number remarks and examples so that cross-references resolve to them.
% alex.sty defines remark as unnumbered and examples on one book-wide counter;
% both are adjusted here (book only) without touching the shared alex.sty.
\newtheorem{remarkn}{Remark}[section]
\crefname{remarkn}{remark}{remarks}
\Crefname{remarkn}{Remark}{Remarks}
\renewenvironment{remark}
  {\pushQED{\qed}\renewcommand{\qedsymbol}{$\triangle$}\remarkn}
  {\popQED\endremarkn}
\counterwithin{examplecounter}{chapter}
\crefname{examplecounter}{example}{examples}
\Crefname{examplecounter}{Example}{Examples}

% Bibliography
\usepackage[backend=biber,style=alphabetic,maxbibnames=99,giveninits=false,
  doi=true,url=true,isbn=false,eprint=true]{biblatex}
\addbibresource{references.bib}
% Let long URLs and DOIs break at any character in the bibliography.
\setcounter{biburlnumpenalty}{100}
\setcounter{biburlucpenalty}{100}
\setcounter{biburllcpenalty}{100}
\renewcommand*{\bibfont}{\small\raggedright}
% Show a URL only when the entry has no DOI.
\DeclareSourcemap{
  \maps[datatype=bibtex]{
    \map[overwrite]{
      \step[fieldsource=doi,final]
      \step[fieldset=url,null]
    }
  }
}

% Better captions and floats
\usepackage{caption}
\usepackage{float}

% Tables for the appendices (notation table spans pages).
\usepackage{booktabs}
\usepackage{longtable}
\usepackage{array}

% Code listings (alex.sty also loads listings; this just adds our default settings)
\lstset{
  basicstyle=\ttfamily\small,
  breaklines=true,
  frame=single,
  numbers=left,
  numberstyle=\tiny,
  showstringspaces=false
}

% Useful math packages not in alex.sty
\usepackage{bm}

% Type-name macros for the textbook. Defined here rather than in alex.sty
% because alex.sty is shared with the papers; book-specific notation lives
% here. \BoolType is a separate name from alex.sty's algorithm2e \Bool
% (which remains available inside algorithm2e environments for pseudocode).
\newcommand{\Bernoulli}[1]{\mathrm{Bernoulli}[#1]}
\newcommand{\BoolType}{\mathrm{bool}}

% Approximate-set macros for the textbook chapter 5 onward.
% alex.sty already provides \ASet (line 453: parametric form supporting
% optional sub/superscripts, defaulting to a ^\pm decoration). \Universe
% is not in alex.sty so we add it here for the discrete universe U.
% \providecommand is used so the textbook stays robust if alex.sty adds
% these names later.
\providecommand{\Universe}{U}

% Classification-measure macros for Part II chapter 7.
% alex.sty already provides \PPV (\RV{PPV}) and \NPV (\RV{NPV}) at
% lines 648 and 649, plus \ACC, \FOR, \FDR, \FPR, \FNR, \TNR, \TPR,
% \TP, \FP, \FN, \TN. The following are not in alex.sty, so we add
% them here. The Greek symbol \pi for prevalence is conventional in
% epidemiology and works inline as $\pi$; the \Prevalence macro just
% makes prose references cleaner.
\providecommand{\Precision}{\mathrm{precision}}
\providecommand{\Recall}{\mathrm{recall}}
\providecommand{\Fone}{F_1}
\providecommand{\Prevalence}{\pi}

% Information-theory macros for Part II chapter 8.
% alex.sty already provides \Entropy (line 610: \FunctionFamily{\mathcal{H}}{X}
% renders as \mathcal{H}(X), with an optional conditioning slot). The textbook
% chapter 8 also needs mutual information I(X;Y) and channel capacity C, which
% alex.sty does not provide. \providecommand keeps things safe if alex.sty
% later adds either name.
\providecommand{\MutualInfo}{I}
\providecommand{\Capacity}{C}
